
% avoid magic strings with objetivos específicos
\newcommand\objetivoEspecificoA{Reubicar la lógica de negocio crítica desde el cliente hacia entornos de ejecución del lado del servidor, de modo que se reduzca la superficie de exposición y se habilite la instrumentación de los flujos sensibles para fines de auditoría y trazabilidad.}

\newcommand\objetivoEspecificoB{Refactorizar la arquitectura interna del \textit{frontend} y del \textit{backend} mediante patrones arquitectónicos modulares que aíslen responsabilidades y reduzcan el acoplamiento, habilitando la evolución independiente de cada dominio funcional y de cada motor tecnológico subyacente.}

\newcommand\objetivoEspecificoC{Diseñar e implementar un sistema transversal de observabilidad estructurada que permita diagnosticar el comportamiento del sistema en producción a través de componentes heterogéneos, soportado por infraestructura como código y notificaciones proactivas.}

\newcommand\objetivoEspecificoD{Consolidar prácticas automatizadas de aseguramiento de calidad continuo a lo largo del ciclo de vida del desarrollo, de modo que la velocidad de cambio se sostenga sin degradar la confiabilidad ni la seguridad del sistema.}

% Activities
% - Objective A: reubicación de lógica crítica al servidor
\newcommand{\objectivoEspecificoAActividadadA}{
  Inventario de la lógica sensible expuesta en el cliente y clasificación de cada flujo según su afinidad con el ciclo de vida de la sesión del usuario.
}
\newcommand{\objectivoEspecificoAActividadadB}{
  Diseño e implementación de funciones \textit{serverless} (AWS Lambda) para la lógica de servicios y las integraciones asíncronas.
}
\newcommand{\objectivoEspecificoAActividadadC}{
  Construcción de un \textit{Backend-for-Frontend} sobre Nuxt~3 para encapsular llamadas a APIs externas (Rama Judicial, SAMAI) y centralizar saneamiento, autenticación y observabilidad.
}
\newcommand{\objectivoEspecificoAActividadadD}{
  Verificación de que la lógica crítica deja de ser inspeccionable desde el cliente y de que los flujos quedan instrumentados para auditoría.
}

% - Objective B: refactorización arquitectónica interna
\newcommand{\objectivoEspecificoBActividadadA}{
  Organización del \textit{frontend} en \textit{bounded contexts} bajo principios de \textit{Domain-Driven Design}, materializados mediante Nuxt Layers.
}
\newcommand{\objectivoEspecificoBActividadadB}{
  Aislamiento del tema visual en una capa independiente del dominio del negocio.
}
\newcommand{\objectivoEspecificoBActividadadC}{
  Refactorización de las funciones Lambda hacia el patrón \textit{Handler--Service--Repository} con repositorios definidos por interfaz.
}
\newcommand{\objectivoEspecificoBActividadadD}{
  Validación del desacoplamiento mediante pruebas unitarias del \textit{service} con implementaciones \textit{mock} del repositorio.
}

% - Objective C: observabilidad transversal
\newcommand{\objectivoEspecificoCActividadadA}{
  Diseño de un esquema estandarizado de telemetría estructurada en formato JSON para frontend y backend.
}
\newcommand{\objectivoEspecificoCActividadadB}{
  Implementación de interceptores HTTP en ambas capas que emitan eventos con identificadores de correlación propagados.
}
\newcommand{\objectivoEspecificoCActividadadC}{
  Aprovisionamiento mediante CloudFormation de filtros de métricas, alarmas, dashboards y notificaciones proactivas hacia Telegram.
}
\newcommand{\objectivoEspecificoCActividadadD}{
  Validación del subsistema correlacionando eventos de extremo a extremo en flujos reales (BFF, microservicios, integraciones externas).
}

% - Objective D: aseguramiento de calidad continuo
\newcommand{\objectivoEspecificoDActividadadA}{
  Configuración de \textit{hooks} de Git con Lefthook (\textit{pre-commit}, \textit{commit-msg}, \textit{post-merge}) integrados con ESLint, Prettier y Commitlint.
}
\newcommand{\objectivoEspecificoDActividadadB}{
  Implementación de pipelines de integración continua en GitHub Actions con validación diferencial sobre archivos modificados.
}
\newcommand{\objectivoEspecificoDActividadadC}{
  Construcción de la suite de pruebas \textit{end-to-end} en Playwright sobre la \textit{Firebase Emulator Suite}.
}
\newcommand{\objectivoEspecificoDActividadadD}{
  Construcción de la suite de pruebas unitarias de las reglas de seguridad de Firestore con Vitest y \texttt{@firebase/rules-unit-testing}.
}

% - Objective Document
\newcommand{\objetivoEspecificoDocumentActividadA}{
  Redacción progresiva.
}
\newcommand{\objetivoEspecificoDocumentActividadB}{
  Revisión con tutor.
}
\newcommand{\objetivoEspecificoDocumentActividadC}{
  Preparación para entrega y defensa.
}

\newcommand{\objetivoEspecificoAActividades}{
  \begin{enumerate}[leftmargin=*, topsep=0pt, itemsep=0pt]
    \item \objectivoEspecificoAActividadadA
    \item \objectivoEspecificoAActividadadB
    \item \objectivoEspecificoAActividadadC
    \item \objectivoEspecificoAActividadadD
  \end{enumerate}
}

\newcommand{\objetivoEspecificoBActividades}{
  \begin{enumerate}[leftmargin=*, topsep=0pt, itemsep=0pt]
    \item \objectivoEspecificoBActividadadA
    \item \objectivoEspecificoBActividadadB
    \item \objectivoEspecificoBActividadadC
    \item \objectivoEspecificoBActividadadD
  \end{enumerate}
}

\newcommand{\objetivoEspecificoCActividades}{
  \begin{enumerate}[leftmargin=*, topsep=0pt, itemsep=0pt]
    \item \objectivoEspecificoCActividadadA
    \item \objectivoEspecificoCActividadadB
    \item \objectivoEspecificoCActividadadC
    \item \objectivoEspecificoCActividadadD
  \end{enumerate}
}
\newcommand{\objetivoEspecificoDActividades}{
  \begin{enumerate}[leftmargin=*, topsep=0pt, itemsep=0pt]
    \item \objectivoEspecificoDActividadadA
    \item \objectivoEspecificoDActividadadB
    \item \objectivoEspecificoDActividadadC
    \item \objectivoEspecificoDActividadadD
  \end{enumerate}
}

\newcommand{\objetivoEspecificoDocumentActividades}{
  \begin{enumerate}[leftmargin=*, topsep=0pt, itemsep=0pt]
    \item \objetivoEspecificoDocumentActividadA
    \item \objetivoEspecificoDocumentActividadB
    \item \objetivoEspecificoDocumentActividadC
  \end{enumerate}
}
